\chapter{Test design}
\label{chap:design}

\section{Functional test design}
\subsection{Decision tables and input categories}
For each rule, the input conditions that could change its outcome were
identified. Requirement 6 gives two conditions for \texttt{VIP15}: the
customer must be a VIP and the original merchandise subtotal must be at
least \pounds150.00. Table~\ref{tab:vip15-decisions} shows how these
conditions became test scenarios. The expected outcomes come from the brief.

\begin{table}[htbp]
    \centering
    \caption{Decision table for \texttt{VIP15} eligibility}
    \label{tab:vip15-decisions}
    \small
    \begin{tabular}{@{}lp{0.18\linewidth}p{0.43\linewidth}p{0.21\linewidth}@{}}
        \toprule
        Customer & Original subtotal & Reason for the input & Expected outcome \\
        \midrule
        VIP & \pounds149.99 & One penny below the minimum. & Reject coupon. \\
        VIP & \pounds150.00 & Exactly at the minimum; a product discount
            reduces the subtotal to \pounds135.00. & Deduct \pounds15.00. \\
        VIP & \pounds150.01 & One penny above the minimum. & Deduct \pounds15.00. \\
        Premium & \pounds150.00 & The subtotal qualifies, so customer type
            determines eligibility. & Reject coupon. \\
        \bottomrule
    \end{tabular}
\end{table}

\subsection{Boundary value analysis}
The values \pounds149.99, \pounds150.00, and \pounds150.01 were selected
to check both sides of the threshold and whether the minimum itself qualified. The
\pounds150.00 scenario also checks which subtotal determines eligibility.
After the product discount, only the original subtotal still meets the
minimum. This lets the test distinguish the two calculations.

\section{Unit tests, integration tests, and mocking}
\subsection{Unit tests and real-service integration}
% Describe what the individual class tests establish. Explain what the
% real-service order tests add when discounts and delivery work together.

\subsection{Isolation with Mockito}
% Explain why you mocked the services in OrderServiceTest and what the
% argument check could show about the handoff to delivery.

\section{Coverage-guided structural testing}
The functional decision table describes the behaviour required by the
brief. The decision-to-decision graphs describe the paths in the supplied
Java implementation. After the functional baseline was run, JaCoCo was
used to locate missed decision outcomes, which were then traced through
the code.

For example, the baseline exercised delivery with valid totals but missed
the branch that rejects a negative total. A structural test was added
with a merchandise total of \(-\)\pounds0.01 and otherwise valid inputs.
The test checks that the delivery service throws
\texttt{IllegalArgumentException}.

